\section{Tessuto connettivo: generalità}

Sotto il termine complessivo di \textbf{tessuto connettivo} sono compresi un gran numero di tessuti, di aspetto e costituzione anche diversi tra loro, distribuiti in tutto il corpo, che, contrariamente al tessuto epiteliale, \textbf{non vengono mai a contatto con l'ambiente esterno}. Essi svolgono, a seconda della loro costituzione, una o più funzioni, correlate più o meno direttamente alla \textbf{connessione strutturale o funzionale} dei distretti dell'organismo.

I tessuti connettivi si interpongono fra tessuti di origine diversa, connettendoli tra loro. Sono costituiti da:
\begin{itemize}
  \item \textbf{componente cellulare}: comprende vari tipi di cellule;
  \item \textbf{matrice extracellulare}: componente amorfa detta sostanza fondamentale $\rightarrow$ componente fibrillare.
\end{itemize}

\rubrica{Classificazione}
\begin{center}
\begin{forest} classalbero
[Tessuto connettivo
  [Propriamente detto
    [Fibroso lasso]
    [Fibroso denso]
    [Elastico]
    [Reticolare]
    [Adiposo]
  ]
  [Di sostegno
    [Cartilagineo]
    [Osseo]
  ]
  [Emolinfopoietico]
]
\end{forest}
\end{center}

% ---------------------------------------------------------------------------
\subsection{Matrice extracellulare}

\rubrica{Sostanza fondamentale}
\begin{itemize}
  \item \textit{fase disperdente acquosa}: in cui sono disciolti elettroliti;
  \item \textit{fase dispersa}: costituita da enzimi, glicoproteine e proteoglicani.
\end{itemize}

\rubrica{Componente fibrillare}
\begin{itemize}
  \item fibre collagene;
  \item fibre reticolari;
  \item fibre elastiche.
\end{itemize}

\rubrica{Proteoglicani}
Costituiti da un \textbf{asse proteico principale} sul quale sono inserite una serie di catene polisaccaridiche complesse dette \textbf{glicosamminoglicani} (GAG) (es.\ acido ialuronico). Funzione: formare una matrice gelatinosa, che promuove la coesione cellulare ed è in grado di trattenere l'acqua a livello degli spazi interstiziali.

\rubrica{Glicoproteine}
Funzione di raccordo tra le molecole della matrice extracellulare e le cellule in essa presenti; costituite da una modesta quota di carboidrati.
\begin{itemize}
  \item \textbf{fibronectina}: permette l'adesione delle cellule al substrato;
  \item \textbf{laminina}: presente nelle lamine basali, permette l'adesione con il collagene;
  \item \textbf{nidogeno} (o \textit{entactina}): tipico delle membrane basali, forma un complesso non covalente con la laminina;
  \item \textbf{condronectina}: presente nella cartilagine, media l'attacco dei condrociti al collagene;
  \item \textbf{osteonectina} (o SPARC, \textit{Secreted Protein, Acidic and Rich in Cysteine}): coinvolta in processi di cicatrizzazione e rimodellamento tissutale.
\end{itemize}

% ---------------------------------------------------------------------------
\subsection{Fibre della matrice}

\rubrica{Fibre collagene}
Alta resistenza, ma poco elastiche; diametro 1--12\,$\mu$m. Tipiche del tessuto connettivo dei legamenti articolari e dei tendini muscolari.
\begin{itemize}
  \item \textbf{tropocollagene}: unità proteica lunga 280\,nm e spessa 1,5\,nm; tre catene si avvolgono a spirale a formare una tripla elica, con sequenza (gly-X-Y)$_{333}$, dove X spesso è prolina e Y è spesso idrossiprolina;
  \item \textbf{microfibrilla}: costituita da molecole di tropocollagene che si associano tra loro sfalsate di $\frac{1}{4}$ a dare una bandeggiatura regolare;
  \item l'associazione di più microfibrille costituisce una \textbf{fibrilla};
  \item l'associazione di più fibrille costituisce una \textbf{fibra collagene}.
\end{itemize}

\rubrica{Fibre reticolari}
Sottili fibre ramificate a formare una lassa rete tridimensionale a maglie larghe, detta \textbf{reticolo}. Impalcatura ramificata resistente ma flessibile, costituita da molecole di tropocollagene che tendono ad aggregarsi lateralmente.
\begin{itemize}
  \item formano l'impalcatura degli organi linfatici;
  \item costituiscono un reticolato intorno alle cellule degli organi parenchimatosi (es.\ fegato, rene, ghiandole endocrine).
\end{itemize}

\rubrica{Fibre elastiche}
Costituite da una regione centrale amorfa contenente \textbf{elastina} e una guaina fibrillare periferica formata da microfibrille tubulari di 10\,nm di \textbf{fibrillina}, diversa sia dall'elastina che dal collagene. Ondulate e ramificate, ritornano alla loro lunghezza iniziale dopo stiramento. Tipiche di alcuni legamenti e del sottocute.

% ---------------------------------------------------------------------------
\subsection{Cellule del tessuto connettivo}

Tutte le cellule del connettivo derivano dalla \textbf{cellula mesenchimale}, che può differenziarsi in fibroblasto, condroblasto (→ condrocita), osteoblasto (→ osteocita), adipocita (bruno o comune) e mastocita.

\rubrica{Fibroblasti}
Cellule affusolate, dotate di lunghe e sottili espansioni citoplasmatiche. Producono acido ialuronico e proteine, che interagiscono a formare proteoglicani (che rendono viscosa la sostanza fondamentale) e subunità proteiche che andranno a formare le fibre extracellulari.
\begin{itemize}
  \item \textbf{funzioni}: produzione della componente fibrillare del tessuto connettivo; elaborazione dei complessi macromolecolari della sostanza amorfa (glicoproteine, proteoglicani);
  \item dopo aver svolto la loro funzione, restano imprigionati nella matrice $\rightarrow$ \textbf{fibrociti}: non si dividono più, ma continuano a sintetizzare molecole della matrice extracellulare.
\end{itemize}

\rubrica{Mastociti}
Cooperano con i fibroblasti nel mantenimento di alcune caratteristiche chimico-fisiche della matrice extracellulare. Partecipano ai meccanismi di difesa rilasciando \textbf{eparina} (anticoagulante) ed \textbf{istamina} (vasodilatatore).

\rubrica{Macrofagi}
Dotati di capacità fagocitiche per eliminare patogeni, sostanze estranee o residui cellulari. Superficie cellulare sollevata in numerose pieghe ed estroflessioni. Derivano dai \textbf{monociti}, che vanno incontro a maturazioni progressive che determinano un aumento del volume del citoplasma; la membrana emette prolungamenti in grado di fissare la cellula al substrato. In grado di migrare nel connettivo e di fagocitare materiali vari come batteri, cellule morte e detriti cellulari.

\rubrica{Adipociti}
Cellule voluminose con citoplasma occupato da un'unica enorme gocciola lipidica $\rightarrow$ riserva energetica. Presenti nel tessuto connettivo fibrillare lasso; costituiscono il tessuto adiposo.

\rubrica{Linfociti, plasmacellule e granulociti}
Globuli bianchi migrati dal sangue con compiti di difesa (produzione di anticorpi e fagocitosi).

% ---------------------------------------------------------------------------
\subsection{Caratteristiche}

\begin{itemize}
  \item \textbf{cellularità}: 50--100\% (presenza di matrice costituita da fibre extracellulari e da sostanza fondamentale);
  \item \textbf{polarità}: 0\% (organizzazione strutturale e funzionale orientata, ma non polarizzata);
  \item \textbf{posizione delle cellule}: assenza di membrana basale su cui sono poste tutte le cellule;
  \item \textbf{vascolarizzazione}: 0--100\% (presenza di vasi che nutrono le cellule);
  \item \textbf{innervazione}: 100\%;
  \item \textbf{proliferazione}: 50\% (il connettivo è un tessuto \textit{stabile}).
\end{itemize}

% ---------------------------------------------------------------------------
\subsection{Funzioni}

\begin{enumerate}
  \item costituire l'impalcatura strutturale del corpo;
  \item trasportare fluidi e sostanze in essi disciolti;
  \item proteggere organi ed apparati;
  \item sostenere e collegare tessuti, organi ed apparati;
  \item immagazzinare riserve energetiche;
  \item difendere l'organismo da sostanze ed agenti estranei.
\end{enumerate}
